\section{An actual strict class sector from a controlled ray host}
\label{sec:cubic-genus-class-sector}

This section supplies the initial nonzero strict classes that were
retained as an input in AC-G90 and AC-G91.  It uses the host of
AC-G92-host and an explicit further class field.  In particular, the
ramification permutation module is not merely counted: it is carried
to an actual quotient of a class group in which every prime above the
coefficient prime splits.  The quotient obtained below is noncanonical.
No numerical \texttt{Big}, \texttt{EC}, or \texttt{Frame} inequality,
and no full source \texttt{Space}/\texttt{At} transport, is inferred.

For a number field \(B\), let \(S_B\) be its primes above an odd
prime \(p\), and write
\[
 E_{B,S}=\mathcal O_B[1/p]^\times,\qquad
 C_{B,S}=\operatorname{Cl}(\mathcal O_B[1/p])(p).
 \tag{G93.1}
\]
The notation \((p)\) means the \(p\)-primary subgroup.  All ray
moduli below are finite and prime to \(p\).  The \(S\)-ray class
group is the ordinary ray class group with the ideal classes of the
primes in \(S_B\) made trivial.  Its class field is therefore split
completely at those primes.

\begin{theorem}[AC-G93-sector: a regular nontrivial class quotient]
\label{AC-G93-sector}
Let \(M/F\) be finite Galois with group \(\Delta\) of order prime
to \(p\).  Choose distinct finite places \(v_1,\ldots,v_r\nmid p\)
of \(F\), each split completely in \(M\), such that every element
of \(E_{M,S}\) has \(p\)-th-power residue at every \(w\mid v_i\).
Suppose \(p\mid(\operatorname N v_i-1)\).

Let \(T/F\) be cyclic of degree \(p^s\), unramified outside the
\(v_i\), split completely at \(S_F\), and with inertia of order
exactly \(p\) at each \(v_i\).  Put \(N=MT\).  The degree
\(p^s\) is not required to equal \(p\).
For every finite field \(k\) of characteristic \(p\), there is a
noncanonical \(k[\Delta]\)-equivariant quotient
\[
 C_{N,S}/p\otimes_{\mathbb F_p} k
       \longtwoheadrightarrow J_k,
 \qquad
 J_k\oplus k\simeq k[\Delta]^{\,r}.
 \tag{G93.2}
\]
Here the action on the displayed \(k\) is trivial.  More precisely,
\(J_k\) is the kernel of a map
\[
 k[\Delta]^r\longrightarrow k,
 \qquad (x_i)\longmapsto\sum_i c_i\operatorname{aug}(x_i),
 \quad c_i\in\mathbb F_p^\times.
 \tag{G93.3}
\]

Let \(\Delta_0\subset\Delta\), let \(H=M^{\Delta_0}\), and put
\(H'=HT\).  If \(\alpha:\Delta_0\to k^\times\) is a nontrivial
character, then the classes on \(G_{H'}\) with coefficients
\(k(\alpha)\) that become unramified everywhere over \(N\) and
vanish at every full \(p\)-decomposition group satisfy
\[
 \dim_k H^1_{\mathrm{str},N}(H',k(\alpha))
       \geq r[\Delta:\Delta_0].
 \tag{G93.4}
\]
Equivalently these are the strict classes for the away-kernel
quotient that kills the normal closure of all away-\(p\) inertia
intersected with \(G_N\).  In the actual cubic application
\(F=H^+\), \([\Delta:\Delta_0]=2\), and this gives at least
\(2r\) classes for each nontrivial adjacent character.
\end{theorem}

\begin{proof}
\emph{1. An exact inertia subgroup, including the \(S\)-unit correction.}
Put \(\mathfrak m=\prod_i\prod_{w\mid v_i}w\).  The localized
ray class sequence is
\[
 E_{M,S}\longrightarrow
 J=\prod_{i,w\mid v_i}(\mathcal O_M/w)^\times
 \longrightarrow \operatorname{Cl}_{\mathfrak m,S}(M)
 \longrightarrow\operatorname{Cl}(\mathcal O_M[1/p])
 \longrightarrow0.
 \tag{G93.5}
\]
For completeness, the map from the residue units is obtained by
weak approximation and a principal ideal.  If that ray ideal becomes
a product of the ideals at \(S_M\) times a ray principal ideal,
the quotient of the two generators is an \(S_M\)-unit.  Conversely
each such unit gives this relation.  This proves exactness at \(J\)
after localization; it is not the unlocalized sequence with the
ordinary unit group left unchanged.  Surjectivity on the right is
proved by representing an ideal class away from \(\mathfrak m\).

Use additive notation for these finite abelian groups.  Let
\(J_p=J(p)\), let \(B\) be the image of \(E_{M,S}\) in \(J\),
and let \(A=\operatorname{Cl}_{\mathfrak m,S}(M)(p)\).
The assumed residue condition says
\[
 B\cap J_p\subset pJ_p.
 \tag{G93.6}
\]
Indeed the projection of a \(p\)-th power in a finite abelian
group to its \(p\)-primary subgroup is a multiple of \(p\).
Thus the image \(I\subset A\) of \(J_p\) satisfies
\[
 I=J_p/(B\cap J_p),\qquad
 I/pI\simeq J_p/pJ_p\simeq\mathbb F_p[\Delta]^r.
 \tag{G93.7}
\]
The last isomorphism is equivariant: the primes above each \(v_i\)
form a free transitive \(\Delta\)-set, and their residue fields
are canonically the residue field of \(F\) at \(v_i\).
Choose a generator of its cyclic \(p\)-Sylow once and transport it.

Now set \(A'=A/pI\).  Since \(pI\) is a subgroup of \(I\),
(G93.7) is an actual injection
\[
 \overline I=I/pI\hookrightarrow A'.
 \tag{G93.8}
\]
This does not assert that the order-\(p\) subgroup of a residue
multiplicative group injects into the original ray class group.
Torsion \(S\)-units may meet that subgroup.  The quotient by
\(pI\), and the use of the full \(S\)-unit condition in (G93.6),
are essential.  In particular no assumption
\(v_p(\operatorname N v_i-1)=1\) has been made.

\emph{2. Produce an everywhere unramified extension of the new host.}
Let \(L/M\) be the finite abelian class field with Galois group
\(A'\).  The modulus, the localized ideal subgroup, and \(pI\)
are all \(\Delta\)-stable.  Consequently \(L/F\) is Galois;
conjugation on its abelian subgroup \(A'\) gives the indicated
\(\Delta\)-action, independently of any lift of an element of
\(\Delta\).
The inertia subgroup of \(L/M\) at every \(w\mid v_i\) is the
corresponding order-\(p\) summand of (G93.8).  It is unramified
at all other finite places and split completely at \(S_M\).

Since \([M:F]\) is prime to \(p\), \(M\cap T=F\).  Hence
\[
 \operatorname{Gal}(N/F)=\Delta\times G,
 \qquad G=\operatorname{Gal}(T/F)\simeq C_{p^s}.
 \tag{G93.9}
\]
The extension \(N/M\) has conductor dividing \(\mathfrak m\):
its ramification at the displayed places is tame of degree \(p\),
and it is otherwise unramified and split at \(S_M\).
Its Artin map therefore factors through \(A\).  It kills \(pI\),
because each local inertia image has order \(p\).
There is a surjection \(A'\to G\), so \(N\subset L\).
Each local summand of \(\overline I\) maps isomorphically onto
the unique subgroup \(G[p]\) of order \(p\).

Put \(K_0=\operatorname{Gal}(L/N)=\ker(A'\to G)\).
At a ramified place the relative inertia in \(L/N\) is the
intersection of that local order-\(p\) subgroup with \(K_0\),
which is zero.  The other finite places were already unramified,
and the places above \(p\) were already completely split.
There is no real ramification in an odd-degree extension.
Thus \(L/N\) is everywhere unramified and completely split at
every prime above \(p\).  It is abelian of \(p\)-power degree.
Global reciprocity gives a surjection
\[
 C_{N,S}\longtwoheadrightarrow K_0.
 \tag{G93.10}
\]
This map is \(\Delta\)-equivariant for the natural subgroup
\(\Delta=\operatorname{Gal}(N/T)\).  To see that no splitting
choice is hidden, conjugation on \(K_0\) by two lifts of an
element of \(\Delta\) agrees: their ratio lies in the abelian
group \(A'\).  Naturality of the Artin map gives the same action
on (G93.10).

The action of \(\Delta\) on \(G\) is trivial by (G93.9).
After choosing \(G[p]\simeq\mathbb F_p\), the map
\(\overline I\to G[p]\) is therefore
\(\sum_i c_i\operatorname{aug}\), with each \(c_i\ne0\).
Its kernel \(J_0\) injects into \(K_0[p]\), and
\[
 0\longrightarrow J_0\longrightarrow
 \mathbb F_p[\Delta]^r\longrightarrow\mathbb F_p
 \longrightarrow0
 \tag{G93.11}
\]
is exact.  This is the promised actual unramified class-field
realization of the nontrivial part of the inertia module.

\emph{3. From an order-\(p\) subgroup to an elementary class quotient.}
For a finite \(p\)-primary \(\mathbb Z_p[\Delta]\)-module \(V\),
the two exact sequences
\[
 0\to V[p]\to V\xrightarrow{p}pV\to0,
 \qquad
 0\to pV\to V\to V/pV\to0
 \tag{G93.12}
\]
show that \(V[p]\) and \(V/pV\) have the same class in the
Grothendieck group of finite \(\mathbb Z_p[\Delta]\)-modules.
Every simple object is killed by \(p\).  Since
\(p\nmid|\Delta|\), both displayed \(\mathbb F_p[\Delta]\)-modules
are semisimple, so equality of their simple multiplicities gives
a noncanonical isomorphism \(V[p]\simeq V/pV\).
Apply this to \(V=K_0\), and use a Maschke complement to the
submodule \(J_0\subset K_0[p]\).  This gives a noncanonical
surjection \(K_0/pK_0\twoheadrightarrow J_0\).
Compose with (G93.10), and extend scalars to \(k\).
Sequence (G93.11) splits by the same averaging argument, proving
(G93.2)--(G93.3).  The averaging denominator is only
\(|\Delta|\), which is prime to \(p\).

\emph{4. Descend the actual strict classes, with the character orientation.}
The extension \(N/H'\) has Galois group \(\Delta_0\) and degree
prime to \(p\), and it trivializes \(k(\alpha)\).
Inflation--restriction and global class field theory identify the
strict subgroup in the theorem with
\[
 H^1_{\mathrm{str},N}(H',k(\alpha))
  =\operatorname{Hom}_{k[\Delta_0]}
       \bigl(C_{N,S}/p\otimes_{\mathbb F_p}k,k(\alpha)\bigr).
 \tag{G93.13}
\]
Indeed restriction is an isomorphism onto invariant classes,
because the positive cohomology of the finite group \(\Delta_0\)
vanishes.  Over \(N\), the required homomorphisms kill all finite
inertia and all full decomposition groups above \(p\), giving
exactly \(C_{N,S}\).  A zero test after the finite local
extensions is equivalent to a zero test before them, by
restriction--corestriction and their prime-to-\(p\) degrees.

For a cocycle \(z\) and \(h\in G_N\), its equivariance is
\[
 z(g h g^{-1})=\alpha(g)z(h).
 \tag{G93.14}
\]
Consequently (G93.13) uses the \(\alpha\)-sector of the class
group, not the \(\alpha^{-1}\)-sector.  This is exactly the
convention in AC-G74, equations (G74.3)--(G74.5).
The noncanonical quotient (G93.2) embeds
\(\operatorname{Hom}_{k[\Delta_0]}(J_k,k(\alpha))\) into (G93.13).
As a \(\Delta_0\)-module, \(k[\Delta]\) is the sum of
\([\Delta:\Delta_0]\) copies of its regular module.
Since \(\alpha\ne1\), the extra trivial summand in (G93.2)
has no homomorphism to \(k(\alpha)\).  The dimension is therefore
\(r[\Delta:\Delta_0]\), proving (G93.4).

Finally, let \(\mathcal N\) be the closed normal closure in
\(G_{H'}\) of away-\(p\) inertia intersected with \(G_N\).
The constructed cocycles vanish on its generators.  The coefficient
is trivial there, and (G93.14) shows they vanish on every conjugate;
continuity then shows they vanish on \(\mathcal N\).
They factor through the exact away-kernel quotient of AC-G90.
Their restrictions to every full \(p\)-local group are zero.
If the literal residual splitting field \(M_0\) is smaller than
\(M\), this gives the same degree-one condition.  On
\(I_w\cap G_{M_0T}\) the coefficient is trivial, so the cocycle
is an additive homomorphism.  Its subgroup
\(I_w\cap G_N\) has index prime to \(p\).  Vanishing on that
subgroup forces vanishing on the whole group: a prime-to-\(p\)
power of every element lies in the subgroup, and multiplication
by that power is invertible in \(k\).  Thus adjoining \(\mu_p\)
or another prime-to-\(p\) splitting factor has not weakened the
literal residual-kernel condition.
This verifies the asserted native local conditions without
inflating an \(H^2\) group from an unrelated quotient.
\end{proof}

\paragraph{Actual cubic application and remaining bounds.}
Take the actual host \(H\) of AC-G83/G87, \(F=H^+\), and a
prime-to-\(p\) Galois extension \(M/F\) splitting the full residual
tuple, its two adjacent characters, and \(\mu_p\).
AC-G92-host supplies all hypotheses of the theorem, for any
\(r\geq1\), with \([T:F]\leq p h_{F,S}(p)\).
Here \(\Delta_0=\operatorname{Gal}(M/H)\) has index two in
\(\Delta\).  The adjacent character
\(\alpha=\chi^{-j}\gamma^{-1}\) is nontrivial on the actual
separated cubic locus, and so is its conjugate.
The theorem supplies at least \(2r\) strict classes on \(HT\).
The coefficient-prime completions are unchanged, so the retained
local nontriviality of \(\alpha\) gives \(H^0(k(\alpha))=0\)
there.  AC-G74's cone correction vanishes: these are actual
nonzero classes in the strict relative \(H^1\) used in AC-G90.

The same quotient also supplies the central classes.  Put
\(\omega=\gamma^{-2}\) and let
\(W=k(\omega\eta_K)\) on the real field \(F\).
Complex conjugation has value \(-1\) on \(\eta_K\) and value
\(1\) on \(\gamma^{-2}\), so \(W\) is a nontrivial odd character.
Take \(\Delta_0=\Delta\) and this character in (G93.4).
There are at least \(r\) strict classes on \(T\) with coefficients
\(W\).  Restriction through \(HT/T\), of degree two, identifies
them with the \(c\)-anti part of the strict
\(H^1(HT,k(\omega))\): the twist by \(\eta_K\) changes invariance
to anti-invariance.  Averaging a representative by
\((b-b^c)/2\) makes \(b^c=-b\) literally.
The local primitives may be averaged at the paired places as well.
Thus one construction supplies adjacent and central anti classes
simultaneously.  On the retained local separated locus the central
\(H^0\) also vanishes; even without this identification, the
equivariant strict-cone surjection lifts a nonzero strict class to
a nonzero anti relative class.

This supplies existence, not an upper bound for the whole class
sector.  The subsequent cup-killing argument still uses the exact
AC-G91 hypotheses, and any numerical comparison with the source
frame, its auxiliary constants, or the original counting budget
must be checked separately.  No estimate for the chosen primes
or the discriminant is obtained here.

\begin{corollary}[AC-G93-fill: actual filled Low data on a higher selected host]
\label{AC-G93-fill}
Retain the actual G83/G87 host and the original local and global
separation hypotheses used in AC-G90 and AC-G91, including
\(\operatorname{comp}_p(\psi)^2=1\), with \(p=269\).
Let \(T/F\) be the selected-place host of AC-G92-host, with
\(r\geq1\), and let \(T_n/T\) be its cyclotomic
\(\mathbb Z_p\)-tower.  Put \(\mathcal H_n=HT_n=KT_n\).
On some sufficiently high \(\mathcal H_n\) there are a nonzero
strict adjacent class, a nonzero strict central anti class, and
an actual polarized long-entry filling with all the local
splittings and normalization equations of AC-G90.
They give the polynomial Low path
\[
 x_X=Xx,\qquad y_X=Xy,\qquad
 z_X=X^2z+(1-X)b
 \tag{G93.15}
\]
from a nonzero central cocycle to the filled nonzero adjacent
point.  Here \(X\) is a polynomial variable, not a tower index.

Under the four explicit native support inputs of AC-G89 and
the retained cyclic automorphic base-change theorem, fresh
native ordinary determinant support exists on each
\(\mathcal H_n\), with the finite, torsionfree, dimension,
\(\mathrm{N4}\), \(\mathrm{P1core}\), and \(\mathrm{P3core}\)
conclusions of AC-G89.  This does not assert that the new
Low path lies in that support ideal, nor that it instantiates
every field of the old single-character \(\mathrm{Space}/\mathrm{At}\).
\end{corollary}

\begin{proof}
\emph{The bounded even dual on the new real base.}
The proof of AC-G91 must be checked on \(T\), which need not
be Galois over \(\mathbb Q\); its earlier displayed Galois
host description is not being imposed on \(T\).
Use the character
\[
 \lambda=\eta_K\bar\psi^2\bar\epsilon^{2e+1},\qquad
 W^\vee(1)=k(\lambda).
 \tag{G93.16}
\]
It is even over \(\mathbb Q\).  The completions of \(T\) at
\(p\) are the unchanged unramified completions of \(F\),
because every \(p\)-place splits in \(T/F\).
Thus its inertia character is \(\bar\epsilon^{2e+1}\),
nontrivial since \(268\nmid2e+1\).  It stays nontrivial
at every \(T_n\), as restriction of \(p\)-power index
does not reduce a character image of order prime to \(p\).

Let \(B_\lambda/\mathbb Q\) be the totally real abelian field
cut out by \(\lambda\).  The already checked AC-G91 argument
gives \(\mu(FB_\lambda(\mu_p))=0\): it is obtained by
finite \(p\)-extension ascent from its specified abelian base.
The further extension
\[
 TB_\lambda(\mu_p)/FB_\lambda(\mu_p)
 \tag{G93.17}
\]
is cyclic of \(p\)-power degree, possibly a proper quotient
of \(\operatorname{Gal}(T/F)\).  Apply exactly the same
finite-\(p\) ascent once more.  This proves
\(\mu(TB_\lambda(\mu_p))=0\), without applying
Ferrero--Washington to a cubic splitting field.
Exact finite-layer class control and the totally real
Euler--duality formula of AC-G74/G76 therefore bound
\(H^1\) for \(k(\lambda)\) uniformly along \(T_n\).

Here the group is the actual away-kernel quotient formed
using \(MT_n/T_n\).  Restricting its classes to the smaller
field \(T_nB_\lambda\) is legitimate: on its away inertia
the coefficient is trivial, and vanishing on the further
prime-to-\(p\) subgroup forces additive vanishing on all of it,
as proved in Step 4 above and AC-G90-quotient.
Thus the bound is for the required full dual \(H^1\),
not merely for an unrelated absolute or strict subgroup.

\emph{Kill the obstruction while retaining both classes.}
The cyclotomic tower is unramified away from \(p\), so
AC-G90-quotient makes the successive away quotients actual
open subgroups of index \(p\), with unchanged normal kernel.
The relative complex uses every full \(p\)-local map.
Exact prime-to-\(p\) descent of arithmetic duality pairs its
degree-two \(W\)-group perfectly with the full
\(H^1(k(\lambda))\).  These are precisely the complexes
and trace pairing of AC-G91, now over \(T_n\).
Nontriviality of \(\lambda\) makes restriction on the latter
\(H^1\) injective.  Bounded dimensions make it eventually
an isomorphism.  Since corestriction after restriction is
multiplication by \(p\), hence zero, corestriction itself
is eventually zero.  Adjointness and perfectness then make
restriction on the strict degree-two obstruction zero.

Choose nonzero adjacent and central classes supplied above.
The adjacent character remains nontrivial on every
\(\mathcal H_n\), so inflation--restriction preserves its
nonzero global cohomology class.  Its local zero conditions
persist.  The central class may be viewed on \(T\) with
the odd nontrivial coefficient \(W\); the same argument
preserves it on \(T_n\), and hence preserves its anti
restriction on \(\mathcal H_n\).
The local separated characters retain their images on
\(p\)-power-index subgroups, so the strict-cone identifications
used in AC-G90 remain valid.  Restrict beyond the dual
stabilization index and then one layer further.
The adjacent class is still nonzero and its mixed obstruction
is zero.  AC-G90-fill supplies an actual long entry and
the paired constant local conjugators.  Its affine averaging
imposes polarity on the cochain, not merely its class.
The inertia witness for normalization persists by taking
a \(p\)-power of the original witness if necessary;
its nontrivial prime-to-\(p\) character values are unchanged
in order.  AC-G90-Low now proves (G93.15), with the nonzero
anti representative \(b\) constructed above.

\emph{Fresh native support, without an unchanged-base map.}
The fields \(\mathcal H_n/T_n\) remain CM, unramified at
finite places, and relatively split at the coefficient and
old automorphic ramification places.  Restriction through
\(HT/H\) and then \(\mathcal H_n/HT\) is through cyclic
\(p\)-power steps.  Rank-three absolute irreducibility survives:
a nontrivial self-twist of order \(p\) would, on determinants,
have cube equal to one.  The retained cyclic base-change
theorem therefore supplies the actual cuspidal representation.
Its parallel regular weights, ordinary flags and exact finite
away types restrict.  Choose a weight \(w\geq2\) from AC-G87
to meet the unchanged arithmetic-point gap condition.

Although these fields need not be Galois over \(\mathbb Q\),
their normal closure over the old Galois host \(H\) has
\(p\)-power degree: it is a compositum of conjugates of
the cyclic selected-place extension and a cyclotomic layer.
Consequently its intersection with the prime-to-\(p\) scalar
field \(\Omega\) of AC-G89 is still \(K\).
The same scalar residual element extends by the identity
on this normal closure.  Chebotarev supplies a fresh rational
auxiliary prime completely split there with the prescribed
scalar residual values.  This verifies the same literal
\(\mathrm{V0ok}\) and \(\mathrm{N4}\) argument.
The source \(\mathrm{p\_unr}\) concerns this auxiliary prime,
not ramification at coefficient prime \(269\).

The degree of \(\mathcal H_n\) over \(\mathbb Q\) has only
the prime factors \(2\) and \(269\), so it cannot contain
\(\mu_{269}\), whose degree is divisible by \(67\).
Choose a coefficient DVR containing the representation
and all embeddings of this host, as permitted by the
native \(\mathrm{Coef}\) and \(\mathrm{Large}\) conditions.
The concrete representation gives the literal \(\mathrm{Pair}\)
and arithmetic \(\mathrm{PT}\) data exactly as in AC-G89.
Under its retained \(\mathrm{LevEx}\),
\(\mathrm{ProModularRingExists}\), \(\mathrm{LamEx}\), and
\(\mathrm{ModU}\) inputs, this gives a proper support ideal
\(J_n\) and the stated native consequences, with inertia-weight
dimension \(3[T_n:\mathbb Q]\).
No cyclotomic-character-preserving map from the old base
frame to its higher subgroup is used, and no prescribed
characteristic-\(p\) curve is shown to lie in \(J_n\).
\end{proof}

\paragraph{Primary input and geometric dictionary.}
The ray sequence used above is the localized version of Milne,
\emph{Class Field Theory}, the 6 May 1997 v3.1 primary mirror,
V.1.5 and its complete printed proof, pp.~115--116:
\url{https://mathweb.ucsd.edu/~csorense/teaching/math204B_W19/Milne_CFT.pdf}.
Existence, the subgroup--field correspondence, the conductor
criterion, and the unramified Hilbert class field are V.3.5--3.7,
Remark 3.8 and Example 3.9, pp.~123--124.  Global reciprocity is
retained as an input; its complete upstream proof is not newly
claimed.  The \(S\)-localization, the quotient by \(pI\), all
equivariant module arguments and the local splitting conditions
have been proved here.  AC-G74 supplies the already audited
strict-cone convention, also rederived in (G93.13)--(G93.14).
In the established SGA~1, Expos\'e~V, \S\S4,7 dictionary,
the finite unramified extension \(L/N\) is a finite etale cover;
complete splitting at \(p\) is imposed in its class-field quotient.
No new theorem of SGA~1 is invoked.
